Il presente lavoro ha studiato una pipeline di $r$-adattività dinamica basata su una rete neurale su grafi, un filtro deterministico di sicurezza geometrica e un solutore ALE--FEM nativo C++/deal.II. La metodologia mantiene fissa la topologia della mesh e modifica le coordinate dei nodi interni a partire da feature derivate dallo stato numerico corrente. Il movimento filtrato viene distribuito su finestre temporali e integrato nella formulazione ALE, senza remapping esplicito nella campagna principale.

Le campagne preliminari hanno mostrato che il miglioramento dell'errore d'interpolazione ottenuto dalla GNN statica non si traduce necessariamente in un miglioramento della soluzione PDE. Le diagnostiche sui trasferimenti hanno inoltre evidenziato un compromesso tra accuratezza e costo per gli operatori $P_1$ e $L^2$, mentre la formulazione ALE ha consentito di trattare il movimento della mesh senza un trasferimento separato dello stato. La campagna end-to-end Stage D, condotta su 28 casi held-out, ha riportato una riduzione media appaiata dell'errore $L^2$ integrato dell'8.04\% e dell'errore $L^2$ finale dell'8.26\% rispetto alla mesh fissa; entrambi i miglioramenti sono stati osservati in tutti i casi esaminati. L'errore in seminorma $H^1$ integrato è diminuito in tutti i casi, mentre l'errore in seminorma $H^1$ finale è migliorato in 21 casi su 28.

Il guadagno di accuratezza è associato a un costo wall-clock mediano circa 8.11 volte superiore alla soluzione su mesh fissa e a interventi del filtro di sicurezza nel 75.4\% delle finestre Stage D. La metodologia non dimostra quindi un vantaggio di efficienza computazionale nel protocollo attuale. I risultati sono limitati a un problema transiente bidimensionale di convezione--diffusione con soluzione manufatta, a una risoluzione iniziale e a specifici intervalli parametrici. Il loro significato è quello di una validazione controllata della pipeline trajectory-aware, non di una dimostrazione di superiorità generale delle GNN per l'adattività di mesh.

La direzione di ricerca più rilevante consiste nell'allineare maggiormente il criterio di adattamento alla qualità della soluzione PDE, mantenendo la verificabilità geometrica e la consistenza ALE. La robustezza rispetto a risoluzione e passo temporale, la riduzione del costo inferenziale e l'estensione a problemi differenti costituiscono i passaggi necessari per stabilire la portata e i limiti del framework.
